\chapter{Challenges}
\label{ch:challenges}

The setting of this thesis is a winged drone flying through a forest of obstacles. The drone is evaluated on two quantities: how far it travels and how much energy it spends along the way. Its body is described by a vector of morphological parameters. Its controller is a neural network that reads the drone's sensors and drives its actuators. For this chapter it is enough to keep in mind the nested problem of \eqref{eq:codesign-bilevel}. There, a body $\vect{m}$ and a controller $\vect{\theta}$ are evaluated together by a performance $F(\vect{m}, \vect{\theta})$, which measures the quality of the flight of the body $\vect{m}$ under the controller $\vect{\theta}$. The only way to obtain the value of $F$ is to simulate that flight.

The three sections that follow present three challenges. Each challenge limits what a search over bodies and controllers can find, and each one therefore shapes the method of this thesis. The first section presents the sequential design challenge. When the body is designed by hand and only the controller is optimized, a better body is never even tried. The second section presents the morphology optimization challenge. When one controller is kept fixed and only the body is optimized, the search prefers the bodies that suit that controller over the bodies with the greatest potential. The third section presents the computational feasibility challenge. The direct remedy, which trains a new controller for every body the search visits, costs more than any search can afford. Together, the three challenges define the requirement that the rest of the thesis sets out to meet: a controller that adapts to the body it is given, at a cost small enough to fit inside the body search.

\section{The Sequential Design Challenge}
\label{sec:challenge-sequential}

Usually, the design of a flying robot starts from the body. An engineer designs the body, or a commercial product provides it. A controller is then optimized for that body. This thesis trains the controller by reinforcement learning; other works tune it by hand or compute it by trajectory optimization. With respect to \eqref{eq:codesign-bilevel}, this procedure corresponds to fixing $\vect{m}$ at a chosen body $\vect{m}_0$ and solving only the inner problem, which returns $\vect{\theta}^{\star}(\vect{m}_0)$. The result is the best controller for the body $\vect{m}_0$. The best controller for one body is not the same thing as the best pair of body and controller. Any body $\vect{m}$ with $F(\vect{m}, \vect{\theta}^{\star}(\vect{m})) > F(\vect{m}_0, \vect{\theta}^{\star}(\vect{m}_0))$ would be a better choice than $\vect{m}_0$, and the sequential procedure never tries it.

Such better bodies exist, because the body design constrains what the controller can do. For example, the energy spent per metre of flight has a lower bound. This bound is set by the airframe, by its lift-to-drag ratio and by the mass it carries, and no controller can change it during flight. Likewise, the manoeuvrability of the drone has an upper bound. This bound is set by the control authority, that is, by the forces and moments the actuators can produce, and no controller can exceed it. The task of this thesis rewards two things: the distance travelled through a cluttered forest and the energy efficiency of the flight. A task that rewards both distance and efficiency asks for a trade-off between the two objectives. A single hand-designed body sits at one point of that trade-off, and that point was chosen before the task was specified. Other bodies may be better on distance, better on efficiency, or better on both objectives at once. Co-design studies report exactly this outcome: when every body is given its own controller, the evolved bodies outperform the conventional design they started from \cite{gupta2021embodied, muff2026hexacopters, bergonti2024codesign}, and the conventional design is never the winner.

\emph{The sequential design challenge is that the body is chosen before any search begins, so a better body can never be found, whatever controller is trained on the chosen one.} The thesis must therefore search over bodies as well as over controllers.
% Possible addition: exam results of the standard mydrone with its specialist
% controller, to serve as the sequential-design baseline.

\section{The Morphology Optimization Challenge}
\label{sec:challenge-morphology}

The opposite of the sequential procedure is to fix the controller and search over the space of possible bodies. This is the cheap way to co-design. A single controller $\vect{\theta}_G$ is trained once, general enough to fly any body in the search space. The same controller then scores every candidate body, and the search maximizes $F(\vect{m}, \vect{\theta}_G)$ instead of $F(\vect{m}, \vect{\theta}^{\star}(\vect{m}))$. These two scores are different functions of $\vect{m}$. The difference between them is not a uniform offset or a uniform scaling, which a good optimizer could ignore. The difference can reverse the order of two bodies. Under the shared controller, the body $\vect{m}_1$ scores higher than the body $\vect{m}_2$. Under their own controllers, $\vect{m}_2$ scores higher than $\vect{m}_1$:
\begin{equation}
F(\vect{m}_1, \vect{\theta}_G) > F(\vect{m}_2, \vect{\theta}_G) \qquad \text{and} \qquad F(\vect{m}_1, \vect{\theta}^{\star}(\vect{m}_1)) < F(\vect{m}_2, \vect{\theta}^{\star}(\vect{m}_2)).
\label{eq:ranking-flip}
\end{equation}
In this case the search under $\vect{\theta}_G$ keeps $\vect{m}_1$ and discards $\vect{m}_2$, which is the body that would have been better.

There is a systematic reason to expect \eqref{eq:ranking-flip} to hold often rather than rarely. A generalist controller is trained on a whole distribution of bodies. Its weights must serve every body of that distribution, so on any single body they are a compromise. On any single body, the generalist therefore achieves less than a controller trained on that body alone. This gap is smallest near the centre of the training distribution. The gap grows towards the edges of the distribution, because there the dynamics of the body differ most from the average dynamics that the controller has learned to expect. The edges, however, are where the search wants to go. A body that is better than the average body of the distribution must differ from that average body: in size, in shape or in control authority. The more a body differs from the average body, the more it is penalized by a controller that was never adapted to it. The fixed controller thus does two things at once. It undervalues the bodies that differ most from its training distribution, which are exactly the bodies the search should find. And it rewards the bodies that it already flies well, which are the bodies close to the average one. A search under a fixed generalist is therefore biased towards the neighbourhood of the training distribution of that generalist. No amount of evolutionary pressure removes this bias, because the bias is built into the fitness function itself.

Training a controller after the search does not repair this bias. Suppose the search under $\vect{\theta}_G$ has returned its best body, and a controller is then trained on that body alone. The pair obtained is the best controller for the body that scored highest under the generalist. This pair is not the best pair either. Consider the body $\vect{m}_2$ of \eqref{eq:ranking-flip}: it scored lower than $\vect{m}_1$ under $\vect{\theta}_G$, but it gains more from a controller of its own. The search discarded $\vect{m}_2$, so no training afterwards can recover it. Specialization applied after the search cannot correct a search that did not account for specialization.

\emph{The morphology optimization challenge is that a search under one fixed controller does not find the bodies with the greatest potential, but the bodies that suit that controller.} What the search needs is a score that reflects what each body can do with a controller adapted to it. The next section shows why the direct way of obtaining such a score is too expensive.

\section{The Computational Feasibility Challenge}
\label{sec:challenge-feasibility}

The previous sections suggest that the co-design of body and controller is a promising direction for this thesis. The direct way to co-design is the one adopted by the co-design works that train their controllers by reinforcement learning \cite{gupta2021embodied, muff2026hexacopters}. These works solve the inner problem exactly: they train a controller from scratch for every body the search evaluates. The cost of this approach in the setting of this thesis can be estimated from three measured quantities. The first is the time needed to train one controller on many bodies. The second is the time needed to train one controller on a single body. The third is the time taken by one search over bodies in which every body is scored by flights alone, with no training inside the loop. \autoref{tab:feasibility} gives the figures.

\begin{table}[htbp]
\centering
\begin{tabularx}{\textwidth}{>{\raggedright\arraybackslash}X r}
\hline
Quantity & Value \\
\hline
Training of one generalist controller (1\,000 PPO iterations, 65\,536 parallel environments) & 14 h \\
Training of one specialist controller on a single body (same settings, on one of the two GPUs) & 1.7 h \\
\hline
Distinct bodies assessed by one body search (64 at first, then 32 new ones at each of 48 updates) & 1\,600 \\
One body search in which bodies are scored by flights alone, with no training inside the loop (Proposed Co-Design Pipeline) & \textbf{94 h} \\
One body search with a controller training per body, at the generalist's cost & 22\,400 h \\
The same, with a specialist trained on the single body (two at a time, one per GPU) & 1\,360 h \\
\hline
\end{tabularx}
\caption[Cost of a controller training inside the body search]{Cost of a controller training inside the body search, from the production runs of this thesis. All times are wall-clock hours of one computing node with two GPUs (NVIDIA V100 PCIe 32GB). The last two rows are hypothetical; the others are measured.}
\label{tab:feasibility}
\end{table}

A single training of the generalist controller takes fourteen hours. A single body search evaluates 1\,600 distinct bodies. Training a controller of that kind for each of them would take two and a half years on the same hardware. A specialist controller, trained on one body alone, is far cheaper: less than two hours on a single GPU. The specialist is the realistic form of the remedy. Even so, with two specialists trained at a time, one per GPU, one training per body brings the search to almost two months. The same search takes four days when every body is scored by flights alone. The factor between the two is almost fifteen at the present size of the search. And the size of the search is not fixed. Every additional body, generation or scenario multiplies the gap again.

The size of the search cannot be kept small. The performance of a body is a random quantity, for two reasons: the forests are generated at random, and the controller acts stochastically. A single flight therefore says little about a body. In this thesis each body is flown on sixty forests before its score is used, and even then the order of the bodies is far from noise-free. Whatever the cost of evaluating one body is, it is paid once per forest. It is paid again every time the controller changes, because the score of a body under an old controller is worthless once the controller has moved. A co-design loop whose inner level is a training run therefore cannot grow along any of the three axes the search needs: scenarios per body, bodies per generation and number of generations. The works that take this route stay under a thousand bodies in total \cite{muff2026hexacopters}, or require over a thousand processors for a single run \cite{gupta2021embodied}. Both works search far smaller design spaces than the one of this thesis, and both use far cheaper evaluations than the flight of a winged drone through a forest.

\emph{The computational feasibility challenge is that training a controller for every body the search visits costs more than any search over bodies can afford.} The feasibility challenge thus rules out the natural answer to the morphology optimization challenge. At the same time, it fixes the shape that any workable answer must have. The controller cannot be trained per body, so it must be trained once, for the whole space of bodies. A controller trained once is, by construction, not adapted to any one body. Its adaptation to the body being scored must therefore come from another mechanism, and that mechanism must cost no more than the flights that score the body anyway. Whether an adaptation that cheap can change which bodies the search finds is the empirical question of this thesis.
